\subsection*{سؤال 16 امتحان 899205 سنة 2008}
\addcontentsline{toc}{subsection}{سؤال 16 امتحان 899205 سنة 2008}

\noindent
\makebox[\textwidth][c]{\includegraphics[width=0.9\paperwidth,keepaspectratio]{ ../../../bagrut_questions/computational_models/DFA_2008_899205_16.png }%
}%

\ifwithsols
\begin{boxSolution}[سؤال 16 امتحان 899205 سنة 2008]
\begin{enumerate}[itemsep=1.5em, label=\alph*.]
    \item .
    \begin{center}
    \includegraphics[width=0.7\paperwidth,keepaspectratio]{../../../images/DFAs/dogs_cats_mice_arabic_version.png}%
\end{center}

\item
لمعرفة القوانين، نلاحظ أن كل الحالات هي حالات قابلة ما عدا الحالة التي في أقصى اليمين ($q_3$ مع أن اسمها لا يظهر في الرسم.) \\
لذلك، المسيرات المسموحة هي كل المسيرات ما عدا تلك التي توصلنا إلى حالة الفخ هذه ($q_3$). \\
هناك طريقتان للوصول إلى هذه الحالة هما:
\begin{enumerate}
\item $q_0 \rightarrow q_1 \rightarrow q_3$ وهذا يحصل عندما نرى \textbf{قطّ فأر}
\item $q_0 \rightarrow q_2 \rightarrow q_3$ وهذا يحصل عندما نرى \textbf{فأر قطّ}
\end{enumerate}

لذلك القواعد التي يصفها الأوتومات:
\begin{itemize}
    \item المسيرة يمكن أن تبدأ بأي حيوان، ويمكن أن تكون فارغة.
    \item لا يمكن أن يظهر فأر وقط متجاورين في المسيرة.
    \item لا توجد قيود على الكلب فيمكنه الظهور في كل مكان في المسيرة.
\end{itemize}
\end{enumerate}
\end{boxSolution}

\clearpage
\fi
